% app_d2_legs.tex -- d = 2 appendix, part 3 of 3 (HT2dPert, 2026-10-16): descendant spacings at eps^2..eps^5, the
% tagged-line test, the eps^3 external-line window W_cs and its tests, the partial eps^4 chi-line window W_cs4chi, and the
% status of every window component with the outlook. Merged from the method write-up of package I.
% Wording of the R_3 cancellation follows the 2026-10-15 review (expected, not established beyond L <= 16).
% Numbers: as in app_d2_eps4.tex (macros where Sec. d2:sec quotes them; otherwise read from the named records).
\section{Two dimensions: descendant spacings and the external-line windows}\label{d2:app:legs}

This appendix supports Section~\ref{d2:desc} and the external-line window of Section~\ref{d2:c5}. The spacings $R_n$ of
Eq.~\eqref{d2:eq:R} test a window at a fixed order without any continuum input; they exposed the defects that remain beyond
the complete $\eps^4$ window, and a tagged external line located the $\eps^3$ defect exactly.

\subsection{Descendant channels}\label{d2:app:desc}
{\small\raggedright\noindent Records: \code{complete\_exact/\allowbreak{}README.md} (\S\S4, 8 and 9),
\code{complete\_exact/\allowbreak{}results/\allowbreak{}\{descendant\_summary.md,\allowbreak{}descendant\_fits.txt\}},
\code{complete\_exact/\allowbreak{}results/\allowbreak{}descendant\_<scheme>\_L<L>.json}, \code{residual\_log/\allowbreak{}results/\allowbreak{}descendant\_e2.txt},
\code{residual\_log/\allowbreak{}results/\allowbreak{}eps5/\allowbreak{}\{descendant\_eps3,\allowbreak{}ce\_wc\_tables\}.txt}, \code{residual\_log/\allowbreak{}fits/\allowbreak{}parent\_check\_wcs/\allowbreak{}verify.txt}; code \code{residual\_log/\allowbreak{}descendant\_e2.py}.\par}

Two kinds of descendants are used: $\chi_{aa}=\partial^a\bar\partial^a\phi$ (one quantum $(a,a)$, level $2a$, primary
$|1\rangle$) and $(\partial\bar\partial)^a\phi^2$ (primary $|2\rangle$). The reference of the latter at $\eps=0$ lies in a
degenerate block and is fixed as the eigenvector of $H_1$ with eigenvalue $E_2^{(1)}$. The small-$L$ solver
\code{complete\_exact/} and its $W_c$ copy \code{residual\_log/\allowbreak{}ce\_wc/} compute these channels exactly in every scheme for
$L\le16$, and $R_0=R_1=0$ exactly everywhere. A spacing counts as restored if it decays like $1/L$.

\paragraph{Order $\eps^2$.}\label{d2:app:desc2} Since $z=0$ for $\chi_{aa}$, the consistent and zero-mode windows coincide on
it. $R_2(\chi_{aa})$ is a sum over two kinds of intermediate states. In the nested ones the external quantum is absorbed and
re-emitted. In the exchange ones the first vertex emits the new $\chi_{aa}$ together with a bubble, and the second absorbs
the old $\chi_{aa}$ with the bubble. The exchange states hold two $(a,a)$ quanta, so the level cutoff leaves the bubble only
$L-4a$ instead of $L$. For the sunset this is a deficit of $\frac2\pi(2a)\times g_1^2\cdot\frac23w_{aa}$, and all other finite
parts cancel, so
\begin{equation}
 \lim_{L\to\infty}R_2(\chi_{aa})=\frac{8a\,w_{aa}}{27\pi}\neq0
\end{equation}
in the consistent, literal and zero-mode schemes; the raw truncation also lacks the vacuum window, which by
Eq.~\eqref{d2:eq:Ev2} grows linearly in $L$. The mass window repairs it. On the one-quantum state
$\langle\chi_{aa}|M'|\chi_{aa}\rangle=2w_{aa}$ and $\ell=2a$, while $\langle1|D^M|1\rangle=0$, so at first order $D^M$ shifts $R_2$ by
exactly $-\frac1{27\pi}\cdot4a\cdot2w_{aa}=-\frac{8a\,w_{aa}}{27\pi}$.

\begin{table}[ht]
\centering\footnotesize\renewcommand{\arraystretch}{1.15}
\begin{tabular}{@{}cllllll@{}}
\toprule
$a$ & $w_{aa}$ & $\frac{8a\,w_{aa}}{27\pi}$ & $R_2(2000)$ & $R_2(20000)$ & with $D^M$, $L=2000$ & with $D^M$, $L=20000$ \\\midrule
1 & 0.25 & 0.023578510 & 0.023584405 & 0.023579100 & $5.90\times10^{-6}$ & $5.90\times10^{-7}$ \\
2 & 0.140625 & 0.026525824 & 0.026532463 & 0.026526487 & $6.64\times10^{-6}$ & $6.63\times10^{-7}$ \\
3 & 0.0976563 & 0.027631067 & 0.027637989 & 0.027631757 & $6.92\times10^{-6}$ & $6.91\times10^{-7}$ \\
5 & 0.0605621 & 0.028559298 & 0.028566467 & 0.028560012 & $7.17\times10^{-6}$ & $7.14\times10^{-7}$ \\
10 & 0.0310454 & 0.029280172 & 0.029287560 & 0.029280905 & $7.39\times10^{-6}$ & $7.33\times10^{-7}$ \\
\bottomrule
\end{tabular}
\caption{The $\eps^2$ spacing of $\chi_{aa}$ in the consistent window at large $L$, and the residual
$R_2-\frac{8a\,w_{aa}}{27\pi}$ left after adding $D^M$ (\code{residual\_log/\allowbreak{}results/\allowbreak{}descendant\_e2.txt}).}
\label{d2:tab:desc2}
\end{table}

Table~\ref{d2:tab:desc2} gives the numbers: with $D^M$ the residual decays like $1/L$, with coefficient $1.2\times10^{-2}$
($a=1$) to $1.5\times10^{-2}$ ($a=10$). The exact small-$L$ solver confirms this for the full schemes at $L\le16$. In the raw
scheme $R_2$ grows ($\chi_{11}$: $0.1077$ at $L=16$). In the consistent and zero-mode schemes $R_2(\chi_{11})$ is constant
($\NTwoRtwoCons$ at $L=16$), and the $\phi^2$ descendant decreases only slowly ($0.1459$ at $L=16$, effective exponent
$-0.27$). In $W$, and in $W_c$, which is identical at this order, $R_2$ is restored: $\NTwoRtwoWcoef/L$ for $\chi_{11}$
($0.000714$ at $L=16$), $0.0146/L$ for $\chi_{22}$ and $0.634/L$ for the $\phi^2$ descendant.

\paragraph{Order $\eps^3$.}\label{d2:app:desc3} None of the schemes $W$, $W_c$ and their variants restores the $\eps^3$
spacing. The external-line window $W_{cs}$ (Appendix~\ref{d2:app:wcs}) reduces it for $\chi_{aa}$ by at least
$\HTpTwoIdentThree\%$ at $L=14$ and 16, but a $1/L$ decay is not established.
Table~\ref{d2:tab:desc3} collects $R_3(\chi_{11})$; the two sources agree where they overlap (scheme $W$ to all printed digits
at $L=6$ to 14).

\begin{table}[ht]
\centering\footnotesize\renewcommand{\arraystretch}{1.15}
\begin{tabular}{@{}lllllll@{}}
\toprule
scheme & $L=4$ & $L=8$ & $L=12$ & $L=14$ & $L=16$ & behavior \\\midrule
raw & 0.1617 & 0.2479 & 0.3355 & 0.3793 & 0.4232 & grows linearly \\
consistent ($=$ zero-mode) & 0.04007 & 0.03941 & 0.04002 & 0.04037 & 0.04071 & constant \\
$\eps^2$ windows only & 0.1874 & 0.2750 & 0.3643 & 0.4089 & & grows linearly \\
$W$ without $D^M$ & 0.05598 & 0.05519 & 0.05577 & 0.05611 & & constant \\
$W$ with $D^{M,(2)}$ only & 0.05812 & 0.06237 & 0.06598 & 0.06749 & & grows slowly \\
$W$ & 0.02316 & 0.02686 & 0.02962 & 0.03074 & 0.03173 & grows, ${\sim}0.009\ln L$ \\
$W{+}V$ & 0.08723 & 0.1138 & 0.1362 & 0.1469 & 0.1573 & grows \\
$W_c$ & 0.01785 & 0.01618 & 0.01579 & 0.01570 & 0.01565 & constant \\
$W_{cs}$ & 0.004225 & 0.000896 & 0.000121 & $-0.000049$ & $-0.000157$ & falls (Appendix~\ref{d2:app:wcstests}) \\
\bottomrule
\end{tabular}
\caption{$R_3$ of $\chi_{11}=\partial\bar\partial\phi$ by scheme, exact (\code{complete\_exact/\allowbreak{}README.md},
\code{residual\_log/\allowbreak{}results/\allowbreak{}eps5/\allowbreak{}\{descendant\_eps3,\allowbreak{}ce\_wc\_tables\}.txt}, \code{residual\_log/\allowbreak{}fits/\allowbreak{}parent\_check\_wcs/\allowbreak{}verify.txt}).
``$\eps^2$ windows only'' is the window of Appendix~\ref{d2:app:window} with no $\eps^3$ part; ``without $D^M$'' and ``with
$D^{M,(2)}$ only'' are $W$ without the mass window and with its $\eps^2$ part only (their $W_c$ versions are identical at this
order).}
\label{d2:tab:desc3}
\end{table}

For the other descendants at $L=16$: $\chi_{22}$ has $R_3=0.04048$ (consistent), $0.03447$ ($W$, growing) and $0.01755$
($W_c$, flat). The $\phi^2$ descendant of level 2 has $0.09469$ (consistent, growing), $0.03864$ ($W$, growing after a minimum of
$0.0074$ near $L=6$ to 8) and $-0.0337$ ($W_c$; $-0.0306$, $-0.0413$ and $-0.0337$ at $L=6$, 10 and 16, not settled).

\paragraph{Structure of the $\eps^3$ defect.} For $a\ll L/2$ the $\eps^3$ spacing of $\chi_{aa}$ is proportional to
$a\,w_{aa}=\frac18\langle\chi_{aa}|\{\hat\ell,M'\}|\chi_{aa}\rangle$ in every scheme. At $L=12$, $R_3/(a\,w_{aa})=0.16$, $0.15$,
$0.15$ for $a=1$, 2, 3 in the consistent scheme, $0.22$, $0.21$, $0.21$ in $W$ without $D^M$, $0.119$, $0.114$, $0.110$ in $W$,
and $0.063$ to $0.068$ in $W_c$. This is the operator structure of a mass window attached to $\chi$ lines, as at $\eps^2$. The
diagnostics locate it:
\begin{itemize}
\item $D^{\rm zm}$ is not the problem at $\eps^3$: $W$ without $D^M$ leaves a constant ($0.056$ for $\chi_{11}$).
\item In $W$, $D^{M,(2)}$ acting at third order produces the slow growth, and $D^{M,(3)}$ shifts the spacing by an almost
 $L$-independent $-0.036$. The growth is the double counting of Appendix~\ref{d2:app:Wc}, which $W_c$ removes.
\item $W_c$ leaves a near constant: at $L=16$, $R_3/(a\,w_{aa})=0.06259$, $0.06242$ and $0.06337$ for $a=1$, 2 and 3,
 changing slowly with $L$. Section~\ref{d2:c5} compares it with the first-order shift of the external-line window,
 $4a\,w_{aa}\,\delta^{(3)}(L)\to-\frac{16}{81\pi}a\,w_{aa}$, with $\frac{16}{81\pi}=\frac23\cdot\frac8{27\pi}=0.06288$.
\end{itemize}
The factor $\frac23$ equals the $\Delta z=\pm2$ share of the angular average of Appendix~\ref{d2:app:mech}, which suggested a
missing channel of the $\eps^3$ exchange sunset that changes the zero-mode number. The tagged-line test of
Appendix~\ref{d2:app:tag} shows instead that the defect is the level budget of the external legs, charged in each slice only
where the legs are present. Whether the $\frac23$ has a simple diagrammatic origin is not derived.

\paragraph{Order $\eps^4$ and higher.} The defect propagates. In $W_c$, $R_4(\chi_{11})=0.019$, $0.037$, $0.049$, $0.057$,
$0.063$, $0.068$, $0.073$ and $0.076$ for $L=2$, 4, \dots, 16; $R_4(\chi_{22})$ rises from $0.034$ ($L=6$) to $0.070$ ($L=16$), and
$R_4$ of the $\phi^2$ descendant from $0.14$ to $0.40$. In $W$, $R_4(\chi_{11})=0.097$ at $L=16$, and orders 4 to 8 are
violated by larger, growing amounts. Appendices~\ref{d2:app:wcstests} and~\ref{d2:app:wcs4chi} give the $\eps^4$ results of
the external-line windows.

\subsection{The tagged-line test}\label{d2:app:tag}
{\small\raggedright\noindent Records: \code{residual\_log/\allowbreak{}results/\allowbreak{}tag/\allowbreak{}tag\_L<L>\_a<a>\_*.json}; code
\code{residual\_log/\allowbreak{}brute\_desc\_tag.py} (exact fraction arithmetic).\par}

In the $\chi_{aa}$ channel the external line is tagged. The initial quantum is a separate copy $\tau_{\rm in}$ of the mode
$(a,a)$ that can only be absorbed, and the final quantum a copy $\tau_{\rm out}$ that can only be created, once. The three
one-quantum states $|\tau_{\rm in}\rangle$, $|\tau_{\rm out}\rangle$ and the regular $|\chi_{aa}\rangle$ all stand for the real
descendant. All three are projected out of every resolvent, as the reference state is in ordinary perturbation theory. The
order-$n$ energy is the sum of the components on $|\tau_{\rm in}\rangle$ and $|\tau_{\rm out}\rangle$ of the order-$n$
right-hand side of the Rayleigh--Schr\"odinger recursion (the vector whose overlap with the reference gives $E^{(n)}$).
Histories that end on the regular $|\chi_{aa}\rangle$ would count those ending on $|\tau_{\rm out}\rangle$ twice and are dropped. The tags count toward the level
budget, and toward $\ell$ in every window, with charges $(\nu_{\rm in},\nu_{\rm out})$.
\begin{itemize}
\item With charges $(2a,2a)$ the construction is identical to the real channel; it reproduces $W_c$ exactly ($L=2$ to 14
 for $a=1$, $L=6$ to 12 for $a=2$).
\item With charges $(0,0)$ the external line is invisible to the cutoff and to the windows (budget-free). Its self-energy
 then has exactly the budget of the $\phi$ self-energy, on which $r$ is tuned.
\end{itemize}

\begin{table}[ht]
\centering\footnotesize\renewcommand{\arraystretch}{1.15}
\begin{tabular}{@{}rllllll@{}}
\toprule
$L$ & charged ($=W_c$): $R_2$ & $R_3$ & budget-free: $R_2$ & $R_3$ & in only: $R_3$ & out only: $R_3$ \\\midrule
2 & 0.004630 & 0.021917 & 0.017361 & 0.028412 & 0.024564 & 0.026879 \\
6 & 0.001808 & 0.016678 & 0.008545 & 0.012127 & 0.012487 & 0.013542 \\
10 & 0.001122 & 0.015932 & 0.005728 & 0.007330 & 0.009882 & 0.010568 \\
14 & 0.000813 & 0.015702 & 0.004317 & 0.005055 & 0.008861 & 0.009368 \\
\bottomrule
\end{tabular}
\caption{$R_2$ and $R_3$ of $\chi_{11}$ with a tagged external line, $W_c$ windows (\code{residual\_log/\allowbreak{}results/\allowbreak{}tag/}).
``In only'' and ``out only'' charge one leg by 2.}
\label{d2:tab:tag}
\end{table}

With the legs budget-free, $R_2\approx0.06/L$ and $R_3\approx0.07/L$ both go to zero (Table~\ref{d2:tab:tag}); charged,
$R_3$ tends to $\NTwoRthreeWcSixteen$. For $a=2$ the budget-free $R_3$ is $0.0313$, $0.0227$, $0.0177$ and $0.0144$ at $L=6$
to 12, against $0.0221$, $0.0195$, $0.0185$ and $0.0180$ in $W_c$. So the whole $\eps^3$ remainder of $W_c$ is the level
budget of the external legs; spectators, nested windows and the auxiliary objects are already treated correctly.

\paragraph{Why average-level charging is exact at $\eps^2$ and not at $\eps^3$.} A mass-type sub-process on a $\chi$ line,
one of the sub-bubbles that $rM'$ cancels, has two external legs, the legs of the $M'$ term. The truncation charges a leg only
in the slices of the sub-process where it is present: an absorbed leg before its absorption, an emitted leg after its
emission. $D^M$ instead charges both legs with their average level $\frac12(\ell_s+\ell_{s'})$ in every slice.
\begin{itemize}
\item At $\eps^2$ the sub-process is the sunset. It has two orderings, normal (no leg present) and exchange (both legs
 present), whose weights are equal at large bubble level, so the average is exact and $R_2\to0$.
\item At $\eps^3$ there are also mixed orderings. In the vertex-corrected sunset (vertices $v_1$, $v_2$, $v_3$; lines $v_1v_2$, twice $v_1v_3$ and twice $v_3v_2$;
 the legs at $v_1$ and $v_2$) the six orderings have $(0,0)$, $(1,0)$, $(0,1)$, $(2,1)$, $(1,2)$ and $(2,2)$ external legs present in
 their two internal slices. Their weights differ, so the response to the level of a leg is not half the response to a
 spectator.
\item The tadpole with a sunset on its loop has both legs at one vertex and charges exactly one leg in every slice, so it is
 unaffected.
\end{itemize}

\subsection{The \texorpdfstring{$\eps^3$}{eps3} external-line window \texorpdfstring{$W_{cs}$}{Wcs}}\label{d2:app:wcs}
{\small\raggedright\noindent Records: \code{residual\_log/\allowbreak{}WINDOW\_SPEC.md} (\S2.5a),
\code{residual\_log/\allowbreak{}large\_L/\allowbreak{}fits/\allowbreak{}delta3\_fits.txt}, \code{residual\_log/\allowbreak{}large\_L/\allowbreak{}results/\allowbreak{}\{validation\_small.txt,\allowbreak{}tag3\_L16000.npz,\allowbreak{}tag3nu\_L4000.npz\}},
\code{residual\_log/\allowbreak{}ce\_wc/\allowbreak{}out\_wcs/}; code \code{residual\_log/\allowbreak{}\{brute\_desc\_tag,\allowbreak{}delta3\}.py},
\code{residual\_log/\allowbreak{}large\_L/\allowbreak{}\{tagchan,\allowbreak{}delta3\_fits\}.py}, \code{residual\_log/\allowbreak{}ce\_wc/\allowbreak{}\{ce\_core,\allowbreak{}ce\_solver,\allowbreak{}ce\_run,\allowbreak{}tables\_wcs\}.py}.\par}

\paragraph{Definition.} Scheme $W_{cs}$ is $W_c$ with one more term at $\eps^3$. The full $\eps^3$ window is
\begin{equation}
 D^{(3)}=\tfrac14\big\{\Delta e_1^{(3)}(L,\hat\ell),M'_0\big\}+\tfrac14\big\{\Delta e_1^{(2)}(L,\hat\ell),M'_1\big\}+\frac{\delta^{(3)}(L)}2\,X_0,
 \qquad \Delta e_1(L,\ell)=e_1(L)-e_1(L-\ell),
 \label{d2:eq:wcsfull}
\end{equation}
with $\delta^{(3)}$ of Eq.~\eqref{d2:eq:wcs} and $M'_k$ the $\eps^k$ part of $M'$.
\begin{itemize}
\item $X$ is $M'$ with every term multiplied by $\eta$ times the sum of the levels of its $\chi$ legs. The crossing sign is
 $\eta=+1$ for terms with one creator and one annihilator ($\chi$ hops and $Z\leftrightarrow\chi$ hops) and $\eta=-1$ for pair
 creation and pair annihilation ($\chi\chi$ or $Z\chi$). Zero-mode legs have level 0.
\item $e_1^{\rm tag}(\lambda;\nu_{\rm in},\nu_{\rm out})$ is $E-E_0$ of the $\phi$ channel $|1\rangle$ at sub-cutoff $\lambda$
 whose external zero-mode line is tagged as in Appendix~\ref{d2:app:tag}. The tags have level 0 in $H_0$ but add
 $\nu_{\rm in}$, $\nu_{\rm out}$ to the budget level; $r$ and $g$ are those of the references at $L$. With zero charges the
 tagged channel equals the $\phi$ channel exactly through $\eps^3$. $e_1^{(3)}$ and $e_1^{{\rm tag},(3)}$ are taken with $r_3$
 removed.
\item $\delta^{(3)}(L)$ is one number per outer cutoff, used in every channel, including the auxiliary ones (rule $W_c$).
 Everything else ($D^{(2)}$, $D^{\rm zm}$ at all orders, $D^{M,(k)}$ for $k\ge4$, rule $W_c$) is unchanged.
\end{itemize}

\paragraph{Derivation of the operator form.} For a hop $s\to s'$ with spectator level $\ell_{\rm spec}$ and $\chi$ legs of
levels $\nu_{\rm in}$, $\nu_{\rm out}$, the slice-resolved window is $\frac12[e_1(L)-e_1^{\rm tag}(L-\ell_{\rm spec};\nu_{\rm
in},\nu_{\rm out})]M'_{s's}$. To first order in the charges this is $\frac12[e_1'\ell_{\rm spec}+\mu(\nu_{\rm in}+\nu_{\rm
out})]M'_{s's}$, with $e_1'=de_1/d\lambda$ and $\mu=\frac12e_1'+\delta$ the response to one unit of leg level; $\delta$ is its
departure from the average-level rule of $D^M$. As $L\to\infty$ in and out legs respond equally, by time reversal; at finite
$L$ they differ by $\mathcal O(1/L)$ (below). Crossing a leg exchanges the slices in which it is present, so a pair leg
responds with $e_1'-\mu=\frac12e_1'-\delta$. With $\frac12(\ell_s+\ell_{s'})=\ell_{\rm spec}+\frac12(\nu_{\rm in}+\nu_{\rm out})$ this is $D^M$
plus $\frac\delta2X$. $\delta$ is calibrated on the smallest charges, $(2,2)$: with
$\delta^{(3)}_\nu(L)=\frac1{2\nu}\big[e_1^{(3)}(L-\nu)-e_1^{{\rm tag},(3)}(L;\nu,\nu)\big]$ for symmetric charges $(\nu,\nu)$,
$\delta^{(3)}=\delta^{(3)}_2$. For $\chi_{11}$ the window is therefore exactly
the slice-resolved one; for larger $a$ it is its linearization, which becomes exact as $L\to\infty$ (below).

\paragraph{What is fixed and what is assumed.} Nothing is fitted: $\delta^{(3)}(L)$ is an exact rational computed from the
tagged $\phi$ channel at each $L$. The descendant tests do not probe the sign of the pair terms, because $R_3$ of any
descendant involves only diagonal hop terms of $X$. The pair sign rests on the crossing argument and on the separate
cancellation of $\ln L$ in the Z--$\chi$ and $\chi\chi$ parts of the $C_5$ response $S(L)$ (Appendix~\ref{d2:app:wcstests}).

\paragraph{$\delta^{(3)}(L)$ at large $L$.} A chain-sum version of the tagged $\phi$ channel
(\code{residual\_log/\allowbreak{}large\_L/\allowbreak{}tagchan.py}) evaluates $\delta^{(3)}(L)$ for every even $L\le\NTwoDeltaLmax$ in one run of 11
minutes. It reproduces the exact rationals for $L=2$ to 16 to $\NTwoDeltaRational$ (relative $2.6\times10^{-14}$), including
the values the $W_{cs}$ runs used, and runs with different $L_{\max}$ agree to $10^{-11}$. Table~\ref{d2:tab:delta3} gives
selected values; Fig.~\ref{d2:fig}b shows the curve.

\begin{table}[ht]
\centering\footnotesize\renewcommand{\arraystretch}{1.15}
\begin{tabular}{@{}rccc@{}}
\toprule
$L$ & $\delta^{(3)}(L)$ & $\delta^{(3)}(L)+\frac4{81\pi}$ & $\delta^{(3)}_4(L)$ \\\midrule
6 & $-0.0148018$ & $+9.2\times10^{-4}$ & $-0.014450$ \\
10 & $-0.0155308$ & $+1.9\times10^{-4}$ & $-0.016733$ \\
16 & $-0.0158036$ & $-8.5\times10^{-5}$ & $-0.016729$ \\
34 (minimum) & $\NTwoDeltaMin$ & $-1.9\times10^{-4}$ & \\
100 & $-0.0158494$ & $-1.3\times10^{-4}$ & $-0.015993$ \\
400 & $-0.0157725$ & $-5.3\times10^{-5}$ & $-0.015806$ \\
1000 & $-0.0157458$ & $-2.7\times10^{-5}$ & $-0.015759$ \\
4000 & $-0.0157278$ & $-8.8\times10^{-6}$ & $-0.015731$ \\
16000 & $-0.0157217$ & $-2.7\times10^{-6}$ & \\
\bottomrule
\end{tabular}
\caption{$\delta^{(3)}(L)$: exact rationals for $L\le16$, chain sums beyond (\code{residual\_log/\allowbreak{}large\_L/\allowbreak{}fits/\allowbreak{}delta3\_fits.txt}).
Third column: distance from $-\frac4{81\pi}=-0.0157190$. Last column: the coefficient $\delta^{(3)}_4$ calibrated on charges $(4,4)$ (the
exact form for $\chi_{22}$).}
\label{d2:tab:delta3}
\end{table}

$\delta^{(3)}(L)$ falls to its minimum at $L=\NTwoDeltaMinL$ and then rises toward the limit of Eq.~\eqref{d2:eq:delta}; the
next term is $+0.013856/L$, followed by $\mathcal O(\ln L/L^2)$. Free-limit fits with a $\ln L/L$ term ($L_{\min}=200$ to
2000, data to $L=16000$) give $\delta_\infty+\frac4{81\pi}$ between $-3\times10^{-11}$ and $+7\times10^{-11}$; without $\ln
L/L$ they miss by $-5\times10^{-7}$ to $-1.5\times10^{-6}$. With the limit fixed, the $\ln L/L$ coefficient is $-0.0058946$,
equal to $-\frac1{54\pi}$. A fit to $L\le4000$ predicts $L=16000$ to $4\times10^{-10}$ with the $\ln L/L$ term
($-2\times10^{-6}$ without). Also $\delta^{(3)}(2)=-0.0093879$, $\delta^{(3)}(4)=-0.0136276$, and $\delta^{(3)}(16)$ is the
rational $-17345700448126303825/1097581272385718321152$. Because $L\le16$ lies on the falling branch before the minimum,
extrapolations in $1/L^2$ from $L\le16$, which do not see the $\ln L/L$ term, miss the limit: they give $-0.0158$ to $-0.0165$
(\code{residual\_log/\allowbreak{}fits/\allowbreak{}parent\_check\_wcs/\allowbreak{}fits.txt}). With the limit identified numerically,
$4|\delta_\infty|/\frac8{27\pi}=\frac23$: the leg coefficient is two thirds of the $\eps^2$ exchange-sunset defect.

\paragraph{The nonlinearities are finite-$L$ effects.} Define the single-leg coefficients $d_{\rm in}$ and $d_{\rm out}$ as
$\delta^{(3)}_2$ with only one leg charged (by 2); with the $(2,2)$ cross term they give $\delta^{(3)}_2=\frac12(d_{\rm in}+d_{\rm
out})+\frac14(\text{cross term})$. At $L\le16$, $d_{\rm in}$ and $d_{\rm out}$ are smaller in magnitude than $\delta^{(3)}$, and
the coefficients for larger charges are larger: at $L=16$, $\delta^{(3)}_4$ and $\delta^{(3)}_6$ (the exact $\chi_{22}$ and $\chi_{33}$ forms)
exceed $\delta^{(3)}_2$ in magnitude by 6\% and 13\% (evaluated from the $\nu=4$ and $\nu=6$ lines of
\code{delta3\_fits.txt}). At large $L$ both effects disappear
(\code{delta3\_fits.txt}):
\begin{itemize}
\item $d_{\rm in}$ and $d_{\rm out}$ both tend to $-\frac4{81\pi}$. Their difference falls like $1/L$ ($1.6\times10^{-5}$ at
 $L=1000$, $3.9\times10^{-6}$ at 4000), so time reversal is restored, and the $(2,2)$ cross term vanishes
 ($-4.0\times10^{-4}$ at $L=1000$, $-1.2\times10^{-4}$ at 4000).
\item $\delta^{(3)}_\nu$ for $\nu=4$, 6, 8 and 16 tends to the same limit, with differences of about $-(\nu-2)\times0.0062/L$ from
 $\delta^{(3)}_2$ (last column of Table~\ref{d2:tab:delta3}).
\end{itemize}
So the linearized window~\eqref{d2:eq:wcsfull}, calibrated on the charges $(2,2)$, is asymptotically exact for every $a$.

\subsection{Tests of \texorpdfstring{$W_{cs}$}{Wcs}}\label{d2:app:wcstests}
{\small\raggedright\noindent Records: \code{residual\_log/\allowbreak{}fits/\allowbreak{}parent\_check\_wcs/\allowbreak{}\{verify,\allowbreak{}fits,\allowbreak{}r3zero,\allowbreak{}blind16\}.txt},
\code{residual\_log/\allowbreak{}results/\allowbreak{}legw/\allowbreak{}legw\_L1000.npz}, \code{residual\_log/\allowbreak{}large\_L/\allowbreak{}fits/\allowbreak{}legw\_logs.txt},
\code{residual\_log/\allowbreak{}ce\_wc/\allowbreak{}\{out\_wcs,\allowbreak{}out\_local,\allowbreak{}out\_wcs4\}/}, \code{residual\_log/\allowbreak{}large\_L/\allowbreak{}results/\allowbreak{}tag3\_L16000.npz}; code \code{residual\_log/\allowbreak{}\{nm\_legw,\allowbreak{}brute\_legw\}.py}.\par}

\paragraph{T1 passes.} $C_0$ to $C_4$, $r_1$ to $r_3$ and $g_1$ to $g_3$ of $W_{cs}$ are identical to those of $W_c$ as exact
rationals for $L=2$ to 16 ($r_4$ differs).

\paragraph{T2 passes, with the same caveat.} $C_5^{W_{cs}}(L)-C_5^{W_c}(L)=\delta^{(3)}(L)\,S(L)$ holds to
$1.2\times10^{-16}$ or better for $L=2$ to 16, with $S(L)$ the $C_5$ response to $\frac12X$ inserted at $\eps^3$
(\code{nm\_legw.py}, checked by brute force at $L=2$, 4, 6). $S(L)$ is known to $L=\NTwoShiftLmax$ and converges with no
$\ln L$:
\begin{center}\footnotesize\renewcommand{\arraystretch}{1.15}
\begin{tabular}{@{}lllllllll@{}}
\toprule
$L$ & 10 & 20 & 50 & 100 & 200 & 400 & 1000 & $\infty$ \\\midrule
$S(L)$ & 0.286252 & 0.319024 & 0.344641 & 0.354908 & 0.360610 & 0.363699 & 0.365693 & \NTwoSinf \\
\bottomrule
\end{tabular}
\end{center}
A free $\ln L$ coefficient comes out at $7\times10^{-7}$ or less (\code{fits.txt}, $L_{\min}=100$ to 400; other tail towers give
at most $4\times10^{-6}$, \code{legw\_logs.txt}). By type of $M'$ term the $\ln L$ coefficients of $S$ are $+\frac1{16}$
(Z--$\chi$ hop), $-\frac1{16}$ (Z--$\chi$ pair), $-\frac1{16}$ ($\chi$ hop), $+\frac1{16}+\frac1{27}$ (pair creation) and
$-\frac1{27}$ (pair annihilation). They cancel within the hop sector and within the pair sector separately, so a leg window with one
coefficient per unit leg level for hops and another for pairs gives no $\ln L$ in $C_5$, whatever the two coefficients are;
the large-$L$ check confirms this to $10^{-6}$. This holds for either sign of the pair terms. With the crossing sign the
coefficients also cancel within the Z--$\chi$ sector and within the $\chi\chi$ sector; with an unsigned leg weight these two
sectors would carry $\pm\frac18$ each. Only the signed weight therefore keeps $C_5$ free of $\ln L$ if the Z--$\chi$ and
$\chi\chi$ parts of the window have independent coefficients, and this is the support for the sign. With the exact $\delta^{(3)}(L)$, $C_5^{W_{cs}}-C_5^{W_c}=-0.005478$, $-0.005625$ and $-0.005700$ at $L=50$,
100 and 200 (evaluated from \code{tag3\_L16000.npz} and \code{legw\_L1000.npz}), tending to $\delta_\infty S_\infty=\NTwoDeltaS$ with no $\ln L$. The window shift is small compared with the
effect of a residual $\ln L$ with $|b|\approx\NTwoBbound$, about $\pm\HTpCfiveTwoLog$ (Section~\ref{d2:c5log}).

\paragraph{T3 for $\chi_{aa}$: passes at $L\le16$, with a possible small remainder.} Table~\ref{d2:tab:wcsr3} gives $R_3$ of
$\chi_{11}$, $\chi_{22}$ and $\chi_{33}$ in $W_{cs}$.

\begin{table}[ht]
\centering\footnotesize\renewcommand{\arraystretch}{1.15}
\begin{tabular}{@{}rlll@{}}
\toprule
$L$ & $\chi_{11}$ & $\chi_{22}$ & $\chi_{33}$ \\\midrule
6 & $+0.001876$ [$0.016678$] & $+0.005487$ [$0.022139$] & $+0.029714$ [$0.047060$] \\
8 & $+0.000896$ [$0.016183$] & $+0.002279$ [$0.019477$] & $+0.012071$ [$0.029985$] \\
10 & $+0.000401$ [$0.015932$] & $+0.000987$ [$0.018459$] & $+0.003591$ [$0.021791$] \\
12 & $+0.000121$ [$0.015789$] & $+0.000345$ [$0.017971$] & $+0.001503$ [$0.019864$] \\
14 & $-0.000049$ [$0.015702$] & $-0.000012$ [$0.017707$] & $+0.000554$ [$0.019012$] \\
16 & $-0.000157$ [$0.015647$] & $-0.000225$ [$0.017554$] & $+0.000046$ [$0.018566$] \\
\bottomrule
\end{tabular}
\caption{$R_3$ in $W_{cs}$, exact, with the $W_c$ value in brackets (\code{residual\_log/\allowbreak{}fits/\allowbreak{}parent\_check\_wcs/\allowbreak{}verify.txt};
$L=16$ from \code{residual\_log/\allowbreak{}ce\_wc/\allowbreak{}out\_wcs/\allowbreak{}wcs\_L16.json}).}
\label{d2:tab:wcsr3}
\end{table}

The table is the basis of the statement in Section~\ref{d2:desc} (at most 3\% of the $W_c$ value left at $L=14$ and 16
for $a=1$ to 3). In $W_{cs}$, $R_3(\chi_{11})$ is $0.0019$ at $L=6$ and below $2\times10^{-4}$ from $L=12$
on, against about $0.0157$ in $W_c$; it crosses zero between $L=12$ and 14. Whether it tends to zero was tested by a hold-out test (\code{blind16.txt}, made when the $L=16$ value already existed):
fits to $L\le14$ predict $L=16$.
\begin{itemize}
\item Free-limit models ($a+b/L^2$, $a+b/L^2+c/L^3$, $a+b/L+c/L^2$; $L_{\min}=6$, 8) give a $\chi_{11}$ limit of $-0.0005$ to
 $-0.0007$ and miss $L=16$ by $5\times10^{-6}$ to $2.8\times10^{-5}$.
\item The zero-limit model $b/L+c/L^2+d/L^3$ misses by $9\times10^{-6}$ ($L_{\min}=8$) and $2.1\times10^{-5}$ ($L_{\min}=6$), equally
 well.
\end{itemize}
So for $\chi_{11}$ a zero limit is not excluded, and a small negative remainder of 3 to 5\% of the $W_c$ value is slightly
favored. For $\chi_{22}$ the same free-limit fits give $-0.0014$ to $+0.0001$, up to about 8\% of its $W_c$ value ($0.0176$); for
$\chi_{33}$ they scatter between $-0.008$ and $+0.022$ and decide nothing. The exact statement behind the window
is the budget-free test of Appendix~\ref{d2:app:tag}, where $R_3$ falls roughly like $1/L$ at $L\le14$ ($R_3\approx0.07/L$). The residual at $L\le16$ is consistent with
finite-$L$ effects (the linearization and the slow convergence of $R_3$ in $W_c$), but $R_3$ itself has not been computed
beyond $L=16$. The exact slice-resolved forms for $\chi_{22}$ and $\chi_{33}$ use $e_1^{\rm tag}(L;2a,2a)$ instead of the
linearization. They give $-0.0004$, $-0.0009$, $-0.0012$ ($\chi_{22}$) and $+0.0019$, $-0.0013$, $-0.0022$ ($\chi_{33}$) at
$L=10$, 12 and 14, drifting slowly negative. The linearized window is closer to zero at these $L$, and both forms share the
same large-$L$ limit.

\paragraph{T3 for the $\phi^2$ descendants: not settled.} The level-2 descendant is shifted by exactly $\frac32\delta^{(3)}(L)$
(ratio $1.5000$ at every $L$), the first-order expectation value of $\frac\delta2X$ in that state. $W_{cs}$ is right for this
state only if its $W_c$ spacing tends to $-\frac32\delta_\infty=\frac6{81\pi}=+0.0236$. Fits of the $W_c$ spacing at $L\le16$
give $+0.005$ to $+0.039$ (\code{fits.txt}), so they neither confirm nor exclude it. At $L=16$ the spacings are
$\NTwoPhiDescTwo$ (level 2; $W_c$: $-0.0337$) and $\NTwoPhiDescFour$ (level 4; $W_c$: $-0.0257$), and the level-4 one is still
moving away from zero. Two-particle descendants have large finite-$L$ terms: their $R_2$ is still $0.04$ at $L=16$, against
$0.0007$ for $\chi_{11}$.

\paragraph{T4 ($\eps^4$ spacing) is not restored by $W_{cs}$.}
\begin{itemize}
\item $W_{cs}$ leaves $R_4$ of every $\chi_{aa}$ exactly unchanged. The direct effect of the $\delta^{(3)}$ term on $R_4$
 ($+0.002115$ at $L=6$) is cancelled exactly by the change it makes in $e_1^{(4)}$ of the auxiliary channels, hence in
 $D^{M,(4)}$.
\item $R_4$ still grows: $R_4(\chi_{11})=0.0489$, $0.0570$, $0.0632$, $0.0682$, $0.0726$ and $0.0765$ at $L=6$ to 16, a local
 slope of about $\NTwoRfourSlope$ per unit of $\ln L$ (least-squares slope over $L=8$ to 16).
\item With budget-free legs (tagged, all windows of $W_{cs}$ through $\eps^3$), $R_4(\chi_{11})=-0.0235$, $-0.0164$, $-0.0118$,
 $-0.0086$, $-0.0062$ and $-0.0043$ at $L=4$ to 14, going to zero. So the growth is the same leg defect at $\eps^4$.
\item The direct transcription to $\eps^4$ (diagnostic scheme $W_{cs4}$: $D^{(4)}\mathrel{+}=\frac12[\delta^{(4)}X_0+\delta^{(3)}X_1]$,
 with $\delta^{(4)}$ from the tagged $\phi$ channel) removes the growth but does not track the budget-free values:
 $R_4(\chi_{11})=-0.0165$, $-0.0133$, $-0.0115$, $-0.0108$ and $-0.0107$ at $L=6$ to 14, flattening at $-0.0107$. The
 reason is that the zero-mode proxy fails at $\eps^4$ (Appendix~\ref{d2:app:wcs4chi}): in the tagged $\phi$ channel $D^{\rm zm}$
 acts on the regular zero modes only, so the tag is a spectator of the zero-mode clusters, whereas the real external zero mode
 joins them. At zero charge the proxy already differs from $e_1^{(4)}$, and $\delta^{(4)}(L)$ grows roughly like $\ln L$ ($-0.0802$
 at $L=6$ to $-0.0991$ at $L=14$; Table~\ref{d2:tab:d4chi}).
\end{itemize}

\begin{table}[ht]
\centering\footnotesize\renewcommand{\arraystretch}{1.15}
\begin{tabular}{@{}rlllll@{}}
\toprule
$L$ & & $C_5$ & $C_6$ & $C_7$ & $C_8$ \\\midrule
14 & $W_{cs}$ & $-20.25163$ & $157.8451$ & $-1440.041$ & $14711.12$ \\
14 & $W_c$ & $-20.24684$ & $157.8302$ & $-1439.928$ & $14710.24$ \\
14 & relative change & $-2.4\times10^{-4}$ & $+9.4\times10^{-5}$ & $-7.8\times10^{-5}$ & $+5.9\times10^{-5}$ \\\midrule
16 & $W_{cs}$ & $-21.06383$ & $165.3880$ & $-1521.286$ & $15664.25$ \\
16 & $W_c$ & $-21.05893$ & $165.3722$ & $-1521.165$ & $15663.31$ \\
16 & relative change & $-2.3\times10^{-4}$ & $+9.6\times10^{-5}$ & $-7.9\times10^{-5}$ & $+6.0\times10^{-5}$ \\
\bottomrule
\end{tabular}
\caption{$C_5$ to $C_8$ in $W_{cs}$ and $W_c$ at $L=14$ and 16, exact rationals rounded (\code{residual\_log/\allowbreak{}ce\_wc/\allowbreak{}out\_wcs/},
\code{residual\_log/\allowbreak{}ce\_wc/\allowbreak{}out\_local/}).}
\label{d2:tab:wcsCn}
\end{table}

\paragraph{Higher coefficients.} $W_{cs}$ changes $C_6$ to $C_8$ by about 0.01\% at $L\le16$ (Table~\ref{d2:tab:wcsCn}), much
less than the differences of about 1\% between $W_c$ and $W$. Every auxiliary object still follows rule $W_c$.

\subsection{The \texorpdfstring{$\eps^4$}{eps4} \texorpdfstring{$\chi$}{chi}-line window \texorpdfstring{$W_{cs4\chi}$}{Wcs4chi}}\label{d2:app:wcs4chi}
{\small\raggedright\noindent Records: \code{residual\_log/\allowbreak{}WINDOW\_SPEC.md} (\S2.5b),
\code{residual\_log/\allowbreak{}fits/\allowbreak{}parent\_check\_p5/\allowbreak{}verify\_p5.txt}, \code{residual\_log/\allowbreak{}ce\_wc/\allowbreak{}\{out\_wcs,\allowbreak{}out\_wcs4,\allowbreak{}out\_wcs4c,\allowbreak{}out\_wcs4x,\allowbreak{}x4\}/}, \code{residual\_log/\allowbreak{}results/\allowbreak{}legw/\allowbreak{}legw\_L1000.npz},
\code{residual\_log/\allowbreak{}results/\allowbreak{}tag/\allowbreak{}\{tag4bf\_*,\allowbreak{}tag4\_*,\allowbreak{}spec4\_*\}.json}; code
\code{residual\_log/\allowbreak{}\{brute\_desc\_tag4,\allowbreak{}brute\_desc\_chi4,\allowbreak{}delta4chi,\allowbreak{}p5\_prepare\_x4\}.py}, \code{residual\_log/\allowbreak{}ce\_wc/\allowbreak{}tables\_p5.py},
\code{residual\_log/\allowbreak{}fits/\allowbreak{}parent\_check\_p5/\allowbreak{}verify\_p5.py}.\par}

\paragraph{Why the zero-mode proxy fails at $\eps^4$.} Through $\eps^3$ a $\chi$ line and the zero-mode line have the same
self-energy function up to $\mathcal O(1/L)$. That is why $D^M$ can use $e_1$ and why the leg window can be calibrated on the
tagged $\phi$ channel. At $\eps^4$ they differ: in the $\phi$ channel the external zero mode belongs to the zero-mode sector and
joins the zero-mode clusters that $D^{\rm zm}$ windows, and a $\chi$ line cannot. Both parts of the $\chi$ mass window are then
off.
\begin{itemize}
\item \emph{Leg part.} The transcription $W_{cs4}$ flattens at $R_4(\chi_{11})=-0.0107$, while the budget-free $\chi_{11}$
 line goes to zero.
\item \emph{Uniform (spectator) part.} Define the $\chi$-line function $e_1^\chi(\lambda)=\Sigma_\chi(\lambda)/w_{11}(\eps)$, with
 $\Sigma_\chi(\lambda)=E(\chi_{11}\text{ line};\lambda)-E_0(\lambda)$ from the budget-free tagged $\chi_{11}$ channel at sub-cutoff
 $\lambda$ under rule $W_c$, without a window on its own line, and the spectator function
 $Y_n(L,\ell)=\Delta e_1^{\chi,(n)}(L,\ell)-\Delta e_1^{(n)}(L,\ell)$ (Table~\ref{d2:tab:Y4}). At fixed $\ell$, $Y_2$ and $Y_3$
 vanish like $1/L$, so $D^M$ with $e_1$ is right for $\chi$-line spectators through $\eps^3$. $Y_4(L,2)$ does not vanish: it is
 $0.0829$, $0.0824$ and $0.0817$ at $L=12$, 14 and 16, about $+0.082$. $Y_4(L,4)$ levels off at $0.1426$ at $L=14$ and 16, $Y_4(L,6)$ is still
 growing there, and for hard spectators ($\ell$ near $L$) every $Y_n$ is large.
\end{itemize}

\begin{table}[ht]
\centering\footnotesize\renewcommand{\arraystretch}{1.15}
\begin{tabular}{@{}rcccccc@{}}
\toprule
$L$ & $Y_2(L,2)$ & $Y_3(L,2)$ & $Y_4(L,2)$ & $Y_4(L,4)$ & $Y_4(L,6)$ & $Y_4(L,L)$ \\\midrule
6 & $-0.0114$ & $-0.0599$ & $+0.0624$ & $-0.0567$ & $-2.4772$ & $-2.477$ \\
8 & $-0.0068$ & $-0.0346$ & $+0.0789$ & $+0.0971$ & $-0.1060$ & $-2.925$ \\
10 & $-0.0045$ & $-0.0225$ & $+0.0826$ & $+0.1312$ & $+0.1085$ & $-3.328$ \\
12 & $-0.0032$ & $-0.0158$ & $+0.0829$ & $+0.1405$ & $+0.1600$ & $-3.706$ \\
14 & $-0.0024$ & $-0.0117$ & $+0.0824$ & $+0.1426$ & $+0.1756$ & $-4.067$ \\
16 & $-0.0019$ & $-0.0090$ & $+0.0817$ & $+0.1426$ & $+0.1803$ & $-4.421$ \\
\bottomrule
\end{tabular}
\caption{The spectator function $Y_n(L,\ell)$, exact rationals rounded (\code{residual\_log/\allowbreak{}ce\_wc/\allowbreak{}x4/}).}
\label{d2:tab:Y4}
\end{table}

\paragraph{Definition.} Scheme $W_{cs4\chi}$ is $W_{cs}$ plus, at $\eps^4$,
\begin{align}
 D^{(4)}&\mathrel{+}=\tfrac14\big\{Y_4(L,\hat\ell),M'_{cc,0}\big\}+\tfrac12\big[\delta_{4\chi}(L)\,X_0+\delta^{(3)}(L)\,X_1\big],
 \label{d2:eq:wcs4chi}\\
 \delta_{4\chi}(L)&=R^{\rm bf}_4(L)-R_4^{W_{cs}}(\chi_{11};L)+\delta^{(3)}(L)-\tfrac14Y_4(L,2).\nonumber
\end{align}
\begin{itemize}
\item $M'_{cc}$ is the $\chi\chi$ part of $M'$ (hops and pair terms without a zero-mode leg). The anticommutator uses the same
 per-state rule as $D^M$, and in auxiliary channels $Y_4$ is evaluated at the outer reference (rule $W_c$).
\item $R^{\rm bf}_4(L)$ is the $\eps^4$ spacing of the budget-free $\chi_{11}$ line at cutoff $L$, and $X_1$ is the $\eps^1$
 part of the leg operator $X$.
\item The second term is a calibration on $\chi_{11}$. On that state the window adds $w_0Y_4(L,2)+4[w_0\delta_{4\chi}+w_1
 \delta^{(3)}]$, with $w_{11}=w_0+w_1\eps+\dots$, $w_0=\frac14$, $w_1=-\frac14$; this is fixed so that the charged line has the
 $\eps^4$ self-energy of the budget-free one. Hence $R_4(\chi_{11})=R^{\rm bf}_4(L)$ by construction, and $\chi_{11}$ is not a
 test.
\item Nothing is fitted: $R^{\rm bf}_4$ and $Y_4$ are exact brute-force rationals at each $L\le16$.
\end{itemize}
Two variants serve as comparisons. Variant ``all'' applies $Y_4$ also to the Z--$\chi$ terms of $M'$; whether a
$Z\leftrightarrow\chi$ sub-process uses $e_1$, $e_1^\chi$ or a mixture is not derived, so the $\chi\chi$ form is the definition
and ``all'' brackets the choice. The leg-only variant keeps the uniform part on $e_1$:
$D^{(4)}\mathrel{+}=\frac12[(\delta_{4\chi}+\frac14Y_4(L,2))X_0+\delta^{(3)}X_1]$.

\begin{table}[ht]
\centering\footnotesize\renewcommand{\arraystretch}{1.15}
\begin{tabular}{@{}rccccc@{}}
\toprule
$L$ & $\delta^{(3)}$ & $\delta_{4\chi,{\rm leg}}$ & $Y_4(L,2)$ & $\delta_{4\chi}$ & proxy $\delta^{(4)}$ \\\midrule
6 & $-0.01480$ & $-0.08012$ & $+0.06243$ & $-0.09572$ & $-0.08023$ \\
8 & $-0.01529$ & $-0.08409$ & $+0.07888$ & $-0.10381$ & $-0.08554$ \\
10 & $-0.01553$ & $-0.08729$ & $+0.08258$ & $-0.10794$ & $-0.09023$ \\
12 & $-0.01567$ & $-0.09007$ & $+0.08293$ & $-0.11080$ & $-0.09469$ \\
14 & $-0.01575$ & $-0.09264$ & $+0.08240$ & $-0.11324$ & $-0.09908$ \\
16 & $-0.01580$ & $-0.09509$ & $+0.08172$ & $-0.11552$ & \\
\bottomrule
\end{tabular}
\caption{Coefficients of the $\eps^4$ windows (\code{residual\_log/\allowbreak{}ce\_wc/\allowbreak{}x4/}, \code{residual\_log/\allowbreak{}results/\allowbreak{}tag/}).
$\delta_{4\chi,{\rm leg}}=\delta_{4\chi}+\frac14Y_4(L,2)$ is the coefficient of the leg-only variant; the proxy $\delta^{(4)}$ is that of
$W_{cs4}$.}
\label{d2:tab:d4chi}
\end{table}

\paragraph{Exact invariances.} For the leg-only variant ($L=2$ to 16) and for $W_{cs4\chi}$ and ``all'' ($L=4$ to 16), the following are
identical to $W_{cs}$ as exact rationals (\code{verify\_p5.txt}): $C_0$ to $C_5$; $r_1$ to $r_4$ and $g_1$ to $g_4$; $E_p$
through $\eps^4$; $R_2$ and $R_3$ of every descendant. $r_5$, $g_5$ and $C_6$ change (at $L=16$, $r_5=10.4473$ against
$10.2495$ and $g_5=-2.2890$ against $-2.2681$). $R_5$ of $\chi_{aa}$ is exactly unchanged, as $R_4$ was under the $\eps^3$
window, and it grows: $R_5(\chi_{11})=0.0737$, $0.0839$, $0.0934$, $0.1042$, $0.1158$, $0.1285$ and $0.1420$ at $L=4$ to 16.
That is the $\eps^5$ instance of the same leg defect.

\paragraph{T4: the $\eps^4$ spacing.} $\chi_{11}$ matches the budget-free line by construction ($-0.0043$ at $L=14$, $-0.0028$ at
16, going to zero). The tests are $\chi_{22}$ and $\chi_{33}$ against their own budget-free lines (Table~\ref{d2:tab:r4chi}).

\begin{table}[ht]
\centering\footnotesize\renewcommand{\arraystretch}{1.15}
\begin{tabular}{@{}llcccccc@{}}
\toprule
 & scheme & $L=6$ & 8 & 10 & 12 & 14 & 16 \\\midrule
$\chi_{22}$ & $W_{cs}$ ($=W_c$) & $+0.0341$ & $+0.0398$ & $+0.0485$ & $+0.0564$ & $+0.0634$ & $+0.0697$ \\
 & proxy $W_{cs4}$ & $-0.0340$ & $-0.0335$ & $-0.0297$ & $-0.0267$ & $-0.0245$ & \\
 & leg only & $-0.0338$ & $-0.0319$ & $-0.0264$ & $-0.0215$ & $-0.0172$ & $-0.0136$ \\
 & $W_{cs4\chi}$ & $-0.0594$ & $-0.0404$ & $-0.0312$ & $-0.0250$ & $-0.0204$ & $-0.0165$ \\
 & budget-free & $-0.0497$ & $-0.0405$ & $-0.0338$ & $-0.0288$ & $-0.0248$ & \\\midrule
$\chi_{33}$ & $W_{cs}$ ($=W_c$) & $+0.0894$ & $+0.0715$ & $+0.0448$ & $+0.0474$ & $+0.0544$ & $+0.0619$ \\
 & proxy $W_{cs4}$ & $+0.0220$ & $-0.0013$ & $-0.0331$ & $-0.0354$ & $-0.0334$ & \\
 & leg only & $+0.0221$ & $+0.0004$ & $-0.0296$ & $-0.0300$ & $-0.0259$ & $-0.0212$ \\
 & $W_{cs4\chi}$ & $-0.2381$ & $-0.0331$ & $-0.0432$ & $-0.0387$ & $-0.0329$ & $-0.0275$ \\
 & budget-free & & $-0.0723$ & $-0.0569$ & $-0.0490$ & $-0.0429$ & \\
\bottomrule
\end{tabular}
\caption{$R_4$ of $\chi_{22}$ and $\chi_{33}$ by scheme, exact (\code{residual\_log/\allowbreak{}fits/\allowbreak{}parent\_check\_p5/\allowbreak{}verify\_p5.txt}). The
$\chi\chi$ form ($W_{cs4\chi}$) and the variant ``all'' give identical $R_4$ for every descendant.}
\label{d2:tab:r4chi}
\end{table}

\begin{itemize}
\item $\chi_{22}$: $W_{cs4\chi}$ closes $\NTwoGapTwo\%$ of the gap between $W_{cs}$ and the budget-free line at $L=14$ ($-0.0204$
 against $-0.0248$) and agrees to $10^{-5}$ at $L=8$ ($-0.04044$ against $-0.04045$). The remaining gap, $+0.0026$, $+0.0037$ and $+0.0044$ at $L=10$, 12 and 14,
 is not shrinking: $L$ times the gap grows ($0.026$, $0.044$, $0.062$). A zero limit for $a\ge2$ is therefore not established.
 The gap reflects the dependence of the $\eps^4$ leg response on the charge, a few percent of the leg coefficient
 (about 3.5\% at $L=14$), which the calibration on $\chi_{11}$, with charges $(2,2)$, cannot see. $R_4(\chi_{22})$ itself falls ($LR_4=-0.31$, $-0.30$, $-0.29$, $-0.26$ at $L=10$ to 16); the budget-free
 $\chi_{22}$ line was not computed at $L=16$.
\item $\chi_{33}$: $\NTwoGapThree\%$ of the gap is closed at $L=14$ ($-0.0329$ against $-0.0429$); the gap is $+0.0137$, $+0.0103$
 and $+0.0101$ at $L=10$, 12 and 14.
\item The leg-only variant misses $\chi_{22}$ by $0.0073$ to $0.0076$ and $\chi_{33}$ by $0.017$ to $0.019$ at $L=12$ and 14, so the uniform
 part $Y_4$ is needed. The proxy $W_{cs4}$ happens to land near $\chi_{22}$ at $L=12$ and 14 but misses $\chi_{11}$ ($-0.0064$ at
 $L=14$) and $\chi_{33}$ ($+0.0095$).
\item The $\phi^2$ descendant of level 2 still grows in every scheme: $R_4=0.0404$, $0.1472$, $0.2080$, $0.2459$, $0.2710$ and
 $0.2885$ in $W_{cs4\chi}$ at $L=6$ to 16 ($W_{cs}$: $0.1301$ to $0.3949$). It is not settled, as at $\eps^3$.
\end{itemize}

\begin{table}[ht]
\centering\footnotesize\renewcommand{\arraystretch}{1.15}
\begin{tabular}{@{}rccccc@{}}
\toprule
$L$ & $W_c$ & $W_{cs}$ & leg only & $W_{cs4\chi}$ & ``all'' \\\midrule
12 & 149.30050 & 149.31434 & 149.29008 & 149.28512 & 149.28594 \\
14 & 157.83017 & 157.84508 & 157.81947 & 157.81354 & 157.81431 \\
16 & 165.37223 & 165.38804 & 165.36121 & 165.35460 & 165.35533 \\
\bottomrule
\end{tabular}
\caption{$C_6$ by scheme, exact rationals rounded (\code{residual\_log/\allowbreak{}fits/\allowbreak{}parent\_check\_p5/\allowbreak{}verify\_p5.txt}); the continuum
value is $\NTwoCsixCont$.}
\label{d2:tab:c6}
\end{table}

\paragraph{$C_6$.} The leg-only variant and the proxy $W_{cs4}$ add the same operator, $\frac12[\delta X_0+\delta^{(3)}X_1]$,
with different coefficients $\delta$ of $X_0$. Their $C_6$ values therefore fix the response of $C_6$ to $\frac12X_0$
inserted at $\eps^4$: it equals the $C_5$ response $S(L)$ of Appendix~\ref{d2:app:wcstests} to $10^{-14}$ at $L=6$ to 14
(replay check \code{p2.c6\_identity}). Relative to $W_{cs}$, the leg-only variant thus shifts $C_6$ by
$\delta_{4\chi,{\rm leg}}(L)S(L)+\delta^{(3)}(L)S_1(L)$, where $S_1$ is the $C_6$ response to $\frac12X_1$; the second term is
$+0.0018$ to $+0.0025$ at $L=6$ to 14 ($S_1=-0.12$ to $-0.16$). The shift is $-0.0256$ at $L=14$ and $-0.0268$ at $L=16$,
and the uniform part of $W_{cs4\chi}$ adds $-0.0035$, $-0.0050$, $-0.0059$ and $-0.0066$ at $L=10$ to 16. The total
$W_{cs4\chi}$ shift relative to $W_{cs}$ is $-0.032$ at $L=14$ and $-0.033$ at $L=16$ (Table~\ref{d2:tab:c6}).

\paragraph{The window coefficient grows with $\ln L$.} $\delta_{4\chi,{\rm leg}}=-0.0801$, $-0.0841$, $-0.0873$, $-0.0901$, $-0.0926$ and $-0.0951$ at
$L=6$ to 16, a local slope of $-0.014$ to $-0.018$ per unit of $\ln L$ ($-0.0183$ between 14 and 16), and $\delta_{4\chi}$ runs
from $-0.0957$ to $-0.1155$. This is forced: in $W_{cs}$, $R_4$ grows with a local slope of about $\NTwoRfourSlope$ per unit of $\ln L$, while the
budget-free line goes to zero, so the slope should approach $-\NTwoRfourSlope$ as the budget-free spacing flattens. Since
the $C_6$ response to the $\eps^4$ leg term is $S(L)$, the leg window's $C_6$ shift then tends to about
$-\NTwoRfourSlope\,S_\infty\ln L\approx-0.010\ln L$, plus a uniform part of order 0.01.

\paragraph{A structural prediction, not a measurement.} If this window is complete at $\eps^4$, $C_6$ in $W_c$ and $W_{cs}$
carries the opposite divergence, about $\NTwoCsixLog\ln L$, from the missing $\chi$-line budget, and $W_{cs4\chi}$ removes it.
This parallels $\eps^3$, where the missing constant leg window left $C_5^{W_c}$ about $0.006$ above $C_5^{W_{cs}}$. It cannot be
tested at $L\le16$: there $C_6\approx\NTwoCsixSixteen$ in every scheme of Table~\ref{d2:tab:c6}, against the continuum value
$\NTwoCsixCont$, $C_6(L)$ still rises by about $\NTwoCsixStep$ per step of 2 in $L$, and the external-line windows shift $C_6$
by about $2\times10^{-4}$ of its value. No statement
about $C_6(\infty)$ is possible from $L\le16$ in any scheme.

\paragraph{Cost.} The $\chi$-line objects are exact brute force. The budget-free line at $L=16$ took 6.1~hours and 76~GB, the
$\chi$-line function 7.3~hours and 110~GB; $L=18$ would need about 20~hours and more than 200~GB per process, more than one
node holds. A large-$L$ test needs an $\eps^6$ version of the shift machinery and a diagrammatic tagged $\chi$ channel for
$\delta_{4\chi}(L)$ and $Y_4(L,\ell)$; neither is built.

\subsection{Status of the window components and outlook}\label{d2:app:status}
{\small\raggedright\noindent Records: \code{residual\_log/\allowbreak{}WINDOW\_SPEC.md} (status table, \S\S2.5a and 2.5b),
\code{zeromode\_exact/\allowbreak{}README.md} (\S\S9 and 10), \code{complete\_exact/\allowbreak{}README.md} (\S\S1 and 5),
\code{continuum/\allowbreak{}fits\_summary.txt}.\par}

\begin{table}[ht]
\centering\footnotesize\renewcommand{\arraystretch}{1.15}
\begin{tabular}{@{}p{6.4cm}p{8.6cm}@{}}
\toprule
component & status \\\midrule
$D^{\rm zm}$ at $\eps^2$, Eq.~\eqref{d2:eq:Dzm} & derived (linked clusters); exact agreement with brute force at $L\le12$; with $D^M$, the fitted limit of $C_4^W$ agrees with the continuum (Section~\ref{d2:c4}) \\
$D^M$ at $\eps^2$, Eq.~\eqref{d2:eq:DM}, or in the $e_1$ form & derived at leading order in $1/L$; verified at $\eps^4$ \\
Bloch $D^{\rm zm}$ at $\eps^3$ & supported at $\eps^5$: no linear term (scheme Bloch); the eigenvalue-diagonal form is excluded ($-0.0051\,L$) \\
self-consistent auxiliary windows (scheme $W$) & double counting; $C_5^W$ diverges like $\NTwoCfiveWLin\,L$ \\
rule $W_c$ (auxiliary windows referred to $L$) & provisional, not derived in full; T1 and T2 pass, T3 fails \\
coupling window $D^V$ & invalid (finite $C_4$ shift of about 0.49; growing $\eps^3$ spacing) \\
$\eps^3$ external-line window ($W_{cs}$, Eq.~\eqref{d2:eq:wcsfull}) & provisional, not derived in full; $\delta_\infty=-\frac4{81\pi}$ identified numerically; linearization asymptotically exact; T1 and T2 pass; T3 for $\chi_{aa}$: at least $\HTpTwoIdentThree\%$ removed at $L=14$ and 16 ($a=1$ to 3); extrapolated remainder 3 to 5\% for $\chi_{11}$ (zero not excluded), up to about 8\% for $\chi_{22}$, undetermined for $\chi_{33}$; cancellation as $L\to\infty$ expected but not established; $\phi^2$ descendants not settled \\
$C_5$ residual logarithm & no $\ln L$ term preferred by the fits with the $\ln^3L/L$ tail; $|b|\lesssim\NTwoBbound$ judged from the median of the fits to $L\le\HTpTwoLfive$ (Section~\ref{d2:c5log}); bounding it well below $\NTwoBbound$ would need $L\gtrsim\NTwoLneeded$ \\
$\eps^4$ $\chi$-line window ($W_{cs4\chi}$, Eq.~\eqref{d2:eq:wcs4chi}) & provisional; exact invariances through $C_5$ and $R_3$; $\chi_{11}$ by construction; $\chi_{22}$, $\chi_{33}$ gaps $\NTwoGapTwo\%$, $\NTwoGapThree\%$ closed at $L=14$ ($\chi_{22}$ gap not shrinking); $\phi^2$ not settled; grows like $\ln L$ \\
$\eps^5$ leg window & open: $R_5(\chi_{aa})$ grows \\
coupling window with $\chi$ legs; higher-order windows; completeness & conjectured from power counting; not computed \\
$C_6$ to $C_8$ with a complete window & not available \\
\bottomrule
\end{tabular}
\caption{Status of the window components in $d=2$. Every component was derived after its effect had been seen and with the
continuum values known (Section~\ref{d2:caveats}); none has a fitted parameter.}
\label{d2:tab:status}
\end{table}

\paragraph{Structure at higher orders.} Since $C_n=-\frac13\Delta_p^3\mathcal E_p^{(n)}$ at every order, only degree three or
more in $p$ survives, and at order $n$ the number of surviving $\chi$ lines is at most $\frac12(4n-6)$: five at $\eps^4$, seven at
$\eps^5$, nine at $\eps^6$. The following is conjectured from power counting.
\begin{itemize}
\item Nested $\chi$ sunsets with $\chi$ spectators start at $\eps^5$. $\{\hat\ell,M'\}$ contains them, because $\hat\ell$ counts
 every quantum present.
\item Coupling bubbles with hard $\chi$ spectators give windows of order one from $\eps^5$ on. Their window behaves like
 $\ln\frac L{L-\ell}$, so the leading form is not enough; $D^V$ as specified is not a correct form, and a correct one is open.
\item The zero-mode effective Hamiltonian is needed at all orders as a matrix $h^{(k)}_{z'z}(\lambda)$, not only its
 eigenvalues; the slope of the mass window acquires a $\ln L$ at $\eps^3$ ($r_3\sim L\ln L$); and divergent three-vertex
 sub-bubbles with spectators appear.
\item The double counting of Appendix~\ref{d2:app:Wc} recurs at every order with nested windows; rule $W_c$ removes it by
 construction.
\end{itemize}

\paragraph{Next steps.}
\begin{enumerate}
\item The residual logarithm of $C_5$: bounding $|b|$ well below $\NTwoBbound$ needs exact $C_5(L)$ to $L\gtrsim\NTwoLneeded$.
\item Large-$L$ tests of the $\eps^4$ window: the prediction $C_6^{W_c}\supset\NTwoCsixLog\ln L$ and $C_6(\infty)$ need the
 $\eps^6$ shift machinery and a diagrammatic tagged $\chi$ channel for $\delta_{4\chi}(L)$ and $Y_4(L,\ell)$, the analogue of
 \code{residual\_log/\allowbreak{}large\_L/\allowbreak{}tagchan.py} with an external $\chi$ line. The $\chi_{22}$ gap would then show whether the
 $a$-dependence of the leg response is a finite-$L$ effect, as it was at $\eps^3$.
\item The $\eps^5$ leg window: $R_5(\chi_{aa})$ grows, and the same construction one order higher is needed before $C_7$.
\item The $\phi^2$ descendants at $\eps^3$ and $\eps^4$ need larger $L$ or a treatment of two-particle legs beyond the
 linearization.
\end{enumerate}

\paragraph{Consequences for the production solver.}
\begin{itemize}
\item The production solver computes its stated truncation correctly through $\eps^8$ (Section~\ref{d2:solvers}); no
 implementation change is needed for what it claims to compute.
\item The windows it implements (consistent, literal) are incomplete in $d=2$ from $\eps^4$ on. Its consistent $C_4$ follows
 $\NTwoCfourConsA-\frac{365}{4374\pi}\ln L+\dots$, and its consistent $C_5$ needs about $+0.2\ln L$
 (\code{residual\_log/\allowbreak{}WINDOW\_SPEC.md}): fits without logarithms miss the continuum value by 1.6 to 1.7 (steps~3
 and~4 of \code{residual\_log/\allowbreak{}fits/\allowbreak{}fits\_c5.txt}, $L\le160$). Fits of $C_4$ to $C_8$ from its data at $L\le30$ (the data available when they were made) fall short of the diagrammatic
 values in magnitude, by 6 to 13\% for the medians over the tail towers (\code{continuum/\allowbreak{}fits\_summary.txt}); they are no
 evidence against those values.
\item The windows matter much more at high order. The diagonal zero-mode window alone changes $C_4$ by about $+0.10$ but $C_8$ by
 $+67\%$ at $L=2$, $+23\%$ at $L=14$ and $+18.5\%$ at $L=18$ (\S9 of \code{zeromode\_exact/\allowbreak{}README.md}).
\item A production implementation of the window must drop the assumption that $D$ is diagonal in the occupation basis.
 It needs the Bloch matrices $h_{z'z}(\lambda)$ for every $\lambda\le L$ (in scheme $W$, zero-mode channels $z\le20$ for the references
 through $\eps^8$, $z\le21$ with descendant channels), the Newton coefficients $e_1(\lambda)$, rule $W_c$ inside every
 auxiliary channel, and the partners projected out of the window objects. It needs the non-diagonal $D^M$, which does not
 vanish on the references and enters $r_3$ and $g_3$. The $\eps^3$ external-line window needs the tagged $\phi$ channel that
 gives $\delta^{(3)}(L)$ and the reweighted operator $X$, and the $\eps^4$ window the budget-free and spectator $\chi$-line
 channels.
\item The exact Fock basis grows fast (463\,GiB at $L=34$). A reliable extrapolation needs exact large-$L$ sums of the kind
 used here, or at least the exact window shifts, to fix the $\ln L$ structure.
\end{itemize}
